% ======================================================================
% Эксперимент 3.5.2. Ошибки прогнозов Тейлора в зависимости от h.
% Ответственный: B — Попадюк Евгений.
%
% Что сдать (main.tex: пункт 1545–1546, ожидаемый результат 1566–1567;
% прогнозы — из п. 3.2, файл 32-taylor.tex):
%   * график ошибок прогнозов P1, P2, P3 из п. 3.2.3 от приращения трафика h;
%   * показать, что прогноз перестаёт работать при h -> s;
%   * указать диапазон h, рабочую точку, точное значение T(lambda+h)=1/(s-h).
%
% Параметры варианта 3: рабочая точка из п. 3.1; s — оттуда же.
% Запуск: python experiments.py 3.5.2. Рисунок: figures/3-5-2-taylor-error.pdf.
% ======================================================================

\taskhead{Эксперимент 3.5.2. Ошибки прогнозов Тейлора}

\taskhead{Постановка эксперимента}
\todo{B: описать, какие прогнозы сравниваются (порядки 1–3 из п. 3.2), диапазон приращений $h$ вплоть до $s$, точное значение для сравнения.}

\taskhead{Программный результат и сравнение}
\todo{B: листинг вызова своих функций из solution.py и запуск python experiments.py 3.5.2.}
\taskfigure[0.75]{figures/3-5-2-taylor-error.pdf}{Ошибки прогнозов Тейлора 1–3 порядков от приращения трафика $h$}{fig:3-5-2}
\todo{B: пояснить график: как ведёт себя ошибка каждого порядка при росте $h$ и что происходит при $h\to s$.}

\taskhead{Вывод}
\todo{B: пределы применимости прогноза фиксированного порядка; связь с радиусом сходимости из п. 3.2.}
